% ===================================================================
% FST-IV: A Scale-Invariant Stability Principle from Particles to Cosmology
%         -- An Overview of the Functional Stability Theory Programme
% Target journal: Foundations of Physics / Entropy
% Status: v1 (programmatic overview, structural and reductive contribution)
% Author: Lukas Geiger, Bernau, Germany
% Version history:
%   v0 -- 2026-04 (initial skeleton, structural draft)
%   v1 -- 2026-05-12 (series-hub: components table, building-blocks status,
%                     bibitems with Zenodo DOIs, AI-disclosure L3,
%                     EN/DE parity)
% ===================================================================

\documentclass[11pt,a4paper]{article}

% --- Packages ------------------------------------------------------
\usepackage[utf8]{inputenc}
\usepackage[T1]{fontenc}
\usepackage{amsmath,amssymb,amsthm}
\usepackage{graphicx}
\usepackage{hyperref}
\usepackage{booktabs}
\usepackage[margin=2.4cm]{geometry}
\usepackage{enumitem}
\usepackage{xcolor}

% --- Theorem environments ------------------------------------------
\newtheorem{axiom}{Axiom}
\newtheorem{definition}{Definition}
\newtheorem{lemma}{Lemma}
\newtheorem{proposition}{Proposition}
\newtheorem{conjecture}{Conjecture}
\theoremstyle{remark}
\newtheorem*{remark}{Remark}

% --- Title ---------------------------------------------------------
\title{Functional Stability Theory ---\\
A Programmatic Hub:\\
Universal Convexity-Uniqueness\\
Across Scales\\[0.5em]
\large\textit{From Particles to Cosmology}}

\author{Lukas Geiger\\
\textit{Independent Researcher, Bernau, Germany}\\
ORCID: \url{https://orcid.org/0009-0005-7296-1534}}

\date{May 2026}

\begin{document}
\maketitle

\begin{abstract}
\noindent
Functional Stability Theory (FST) is a research programme that proposes
a single stability principle -- the \emph{Universal Convexity-Uniqueness
criterion} (UCU) -- as a unifying framework for emergent structure
across physical scales, from thermodynamic equilibria of fundamental
parameters (FST-I) through autocatalytic chemistry (FST-II) and
biological evolutionary stability (FST-III) up to cosmological dark
energy and turbulence. This overview synthesises the programme. The
universal game $(x_L, G_L, D_L)$ is summarised together with the
Maximum Entropy Production Principle (MEPP) initial-condition axiom,
and the three-lemma endgame UCU1 (strict convexity), UCU2 (existence
of minimiser) and UCU3 (quantitative quadratic lower bound) is
formulated as the technical core. The renormalisation group (RG)
provides the mathematical bridge between scales. We tabulate the
status of UCU1/UCU2/UCU3 across applications (thermodynamics,
chemistry, biology, K41 turbulence, cosmological dark energy,
Yang--Mills mass gap) and document the methodological lesson learned
from the cosmological route (\emph{Pattern~A}): a positivity
reformulation of an open conjecture is in general only a translation
of language, not progress, unless its proof strength is independently
verified. This article is an \emph{honest survey}: it describes what
the programme delivers, what remains conditional, and which sub-problems
are open (HD3' for Beal, full UCU3 for Yang--Mills, L4 RG-matching
for the dark-energy/CRM-VI line).
\end{abstract}

\textbf{Keywords:} stability, entropy production, renormalisation group,
Nash equilibrium, scale invariance, fine-tuning, dark energy,
turbulence, mass gap.

% ===================================================================
\section{Introduction}
\label{sec:intro}
% ===================================================================

The fine-tuning problem -- the empirical observation that a wide range
of physical parameters (Standard-Model couplings, the cosmological
constant, the chemical reaction networks compatible with life) appears
to be selected with high precision among all logically possible values
-- is usually treated as scale-specific: anthropic at the cosmological
scale, kinetic at the chemical scale, evolutionary at the biological
scale. Functional Stability Theory (FST) explores the alternative
hypothesis that these selection statements are different manifestations
of a single stability principle.

The programme has a fixed shape across scales:
\begin{enumerate}[label=(\roman*),leftmargin=2em]
\item A configuration space $X$ on which a functional $F$ is defined.
\item A conditional system: $F$ admits a unique minimiser provided
      a strict convexity property (UCU1), an existence property (UCU2)
      and a quantitative quadratic lower bound (UCU3) hold on
      a closed convex feasibility set $K \subset X$.
\item A renormalisation-group (RG) bridge connecting two scales,
      ensuring that the minimiser at one scale induces a stability
      statement at another.
\end{enumerate}
The contribution of this overview is therefore not a new theorem,
but the \emph{programmatic} claim that scales as different as the
Standard-Model parameter selection (FST-I) and dark-energy late-time
acceleration (CRM-VI) admit the same conditional structure. Earlier
papers in the series treat individual scales; the present article
maps them onto a single architecture and identifies which sub-problems
remain open.

% ===================================================================
\subsection*{FST Series at a Glance}
\label{subsec:series-glance}
% ===================================================================

This paper serves as the \emph{programmatic hub} for the Functional
Stability Theory series. Table~\ref{tab:series-components} lists the
direct application-scale companion papers (FST-I, FST-II, FST-III,
FST-DE) and the major published FST building blocks in the mathematics
and physics lines that this overview references.

\begin{table}[h]
\centering
\renewcommand{\arraystretch}{1.18}
\small
\begin{tabular}{lllc}
\toprule
\textbf{Paper} & \textbf{Scale / Target} & \textbf{Reference} &
\textbf{Series role} \\
\midrule
\multicolumn{4}{l}{\textit{Direct application companions (FST-PRACTICAL)}} \\
FST-I  & Thermodynamic / SM parameters   & \cite{GeigerFST-I-2026}  & hasPart \\
FST-II & Chemical / autocatalysis        & \cite{GeigerFST-II-2026} & hasPart \\
FST-III& Biological / protein folding    & \cite{GeigerFST-III-2026}& hasPart \\
FST-DE & Cosmological / dark energy      & \cite{GeigerFST-DE-2026} & hasPart \\
\midrule
\multicolumn{4}{l}{\textit{Mathematical building blocks (FST-MATHEMATICS)}} \\
RH v2.1 (Even Dominance)        & spectral stability      & \cite{GeigerRH-EvenDom-2026} & references \\
Zeta Zoo (NE-B Boundary)        & classification          & \cite{GeigerZetaZoo-2026}    & references \\
Zookeeper (CCM route)           & RH technical core       & \cite{GeigerZookeeper-2026}  & references \\
FST Spectrum Duality            & physics-master          & \cite{GeigerSpectrumDuality-2026} & references \\
Hodge Positivity                & AP / algebraicity       & \cite{GeigerHodge-2026}      & references \\
BSD Positivity                  & elliptic curve heights  & \cite{GeigerBSD-2026}        & references \\
P vs NP Entropy                 & complexity              & \cite{GeigerPvsNP-2026}      & references \\
\midrule
\multicolumn{4}{l}{\textit{Physical building blocks (FST-PHYSICS)}} \\
Yang--Mills Mass Gap            & strong-coupling lattice & \cite{GeigerYM-2026}         & references \\
Navier--Stokes Regularity       & 3D NS attractor route   & \cite{GeigerNS-2026}         & references \\
NS-LDI Framework                & TLL/LDI bridge          & \cite{GeigerNSLDI-2026}      & references \\
K41 Turbulence (joint minimiser)& Kolmogorov spectrum     & \cite{GeigerK41-2026}        & references \\
\bottomrule
\end{tabular}
\caption{FST series components. ``hasPart'' (Zenodo Related Identifier
\texttt{hasPart}) marks the four direct application-scale companions
that are co-published with this overview. ``references'' marks
published FST building blocks that the overview cites but does not
formally bundle. All identifiers are persistent Zenodo DOIs; see the
bibliography for resolvers.}
\label{tab:series-components}
\end{table}

% ===================================================================
\section{The Universal Game}
\label{sec:universal-game}
% ===================================================================

\subsection{Two Axioms}

\begin{axiom}[Initial Condition]
The cosmos starts in a state of maximal gradient and minimal structure
(low entropy, high free-energy gradient).
\end{axiom}

\begin{axiom}[Evolution: MEPP]
Open non-equilibrium subsystems evolve, on the timescales of interest,
along trajectories consistent with the Maximum Entropy Production
Principle (MEPP)~\cite{MartyushevSeleznev2006}.
\end{axiom}

\subsection{Scale Triple}

For each physical scale $L$ we associate a triple
\[
\mathcal{S}_L \;=\; (x_L,\; G_L,\; D_L),
\]
where $x_L$ is the configuration variable (field, population fraction,
density profile), $G_L = (S_L, \Pi_L)$ is the local game
(strategy space $S_L$ and payoff $\Pi_L$), and $D_L$ is the dynamics
$\dot{x}_L = D_L(x_L; \Pi_L)$.

\subsection{Stability Definition}

\begin{definition}[Stable State]
$x_L^*$ is \emph{stable} at scale $L$ if it is a locally attractive
fixed point of $D_L$ (Lyapunov sense) and robust against invasion
or perturbation (Nash/ESS criterion at the game level).
\end{definition}

The choice of payoff proxy at each scale (entropy production
$\dot S$, free-energy destruction $-\Delta G$, mean fitness $\bar f$,
etc.) is dictated by the physical setting; the FST claim is that the
\emph{long-time stable} attractor correlates with the maximum payoff
proxy under MEPP-compatible boundary conditions.

% ===================================================================
\section{The Universal Stability Criterion (UCU)}
\label{sec:ucu}
% ===================================================================

The technical core of the programme is a three-lemma endgame
formulated in the convex-analysis language of reflexive Banach spaces.
We refer to \texttt{CORE/UCU\_FORMAL\_DEFINITIONS.md}
for the formal definitions and the standard literature
(Brezis~\cite{Brezis2011}, Ekeland--T\'emam~\cite{EkelandTemam1999},
Evans~\cite{Evans2010}). Throughout this section, $X$ is a reflexive
real Banach space, $K \subset X$ is non-empty, closed and convex, and
$F : K \to \mathbb{R} \cup \{+\infty\}$ satisfies (A1) properness,
(A2) weak lower semicontinuity, (A3) coercivity.

\begin{lemma}[UCU1 -- Strict Convexity]
\label{lem:ucu1}
$F$ is \emph{strictly convex} on $K$:
$F(\lambda x + (1{-}\lambda) y) < \lambda F(x) + (1{-}\lambda)F(y)$
for all $x \neq y$ in $K$ and $\lambda \in (0,1)$.
\end{lemma}

\begin{lemma}[UCU2 -- Existence]
\label{lem:ucu2}
Under (A1)--(A3), $F$ attains its infimum on $K$:
there exists $x^* \in K$ with $F(x^*) = \inf_K F$.
\end{lemma}

\begin{lemma}[UCU3 -- Quantitative Quadratic Lower Bound]
\label{lem:ucu3}
There exists $\mu > 0$ such that for all $x \in K$,
\[
F(x) \;\ge\; F(x^*) \;+\; \tfrac{\mu}{2}\,\|x - x^*\|^2.
\]
UCU3 splits into a \emph{local} and a \emph{global} version
(see Section~\ref{sec:cross-scale}).
\end{lemma}

\noindent\textbf{Consequence (UCU-Theorem, conditional form).}
If UCU1 + UCU2 + UCU3 hold on $K$, then $F$ admits a \emph{unique}
minimiser, and the minimum is \emph{quantitatively isolated} in
$\|\cdot\|$. The first two lemmas are classical (Tonelli direct
method); UCU3 is the application-specific input that ties the
abstract result to the physics of the scale.

\begin{remark}
The role of UCU3 is the same as the role of a \emph{spectral gap}
in the resolvent picture of \cite{GeigerZookeeper-2026}: it converts a
qualitative uniqueness statement into a quantitative stability
estimate. Across scales, UCU3 reads ``no flat direction at the
minimum''.
\end{remark}

% ===================================================================
\section{Renormalisation Group as Mathematical Bridge}
\label{sec:rg-bridge}
% ===================================================================

The RG provides the formal apparatus for relating the universal game
at scale $L$ to that at scale $L+1$. A Kadanoff block transformation
$R_b : (x_L, G_L, D_L) \mapsto (x_{L+1}, G_{L+1}, D_{L+1})$ is required
to preserve at least one of the following invariants:
\begin{enumerate}[label=(\Alph*),leftmargin=2em]
\item the fixed-point structure (number and type of stable attractors),
\item the best-response geometry (sign pattern of the Jacobian
      around $x_L^*$),
\item the existence of a Lyapunov / potential function.
\end{enumerate}
Whenever (A)--(C) is preserved, a UCU-stable scale-$L$ minimiser
induces a UCU-stable scale-$(L{+}1)$ minimiser.

\paragraph{Compact statement.}
The RG-bridge is the mechanism by which the conditional UCU theorem
becomes \emph{transferable}. In practice, the bridge has been worked
out only in fragments (see the open problems in
Section~\ref{sec:open}); for the cosmological RG-matching of the
$R + \gamma R^2$ sector (CRM-VI), the L4 step remains open.

% ===================================================================
\section{Cross-Scale Status Map}
\label{sec:cross-scale}
% ===================================================================

Table~\ref{tab:status} summarises the present status of UCU1, UCU2,
UCU3 (local and global) across the six application scales of the
programme.

\begin{table}[h]
\centering
\renewcommand{\arraystretch}{1.18}
\begin{tabular}{lllllll}
\toprule
\textbf{Scale} & \textbf{Functional} & \textbf{UCU1} & \textbf{UCU2}
& \textbf{UCU3 local} & \textbf{UCU3 global} \\
\midrule
FST-I (Thermo / SM)        & $\Sigma[\alpha,m_e/m_p,\dots]$ &
\checkmark & \checkmark & plateau & open \\
FST-II (Chemistry)         & $G[\text{hypercycle}]$ &
\checkmark & \checkmark & local OK & open \\
FST-III (Biology / fold)   & $F_{\text{fold}}[\phi]$ &
\checkmark & \checkmark & local OK & open \\
K41 (Turbulence)           & $F[E(k)]$ &
\checkmark & \checkmark & \checkmark $(g^*=\tfrac12)$ & LSI open \\
DE / CRM-VI (Cosmology)    & $\Phi(\rho)$ &
\checkmark & \checkmark & \checkmark & L4 RG-match open \\
YM (Mass gap)              & $\mathcal{E}[A]$ &
? & ? & ? & UCU3 = \emph{mass-gap question} \\
\bottomrule
\end{tabular}
\caption{Status of UCU1/UCU2/UCU3 across application scales.
\checkmark{} = established conditional on standard inputs (Brezis,
Ekeland--T\'emam direct method); ``open'' = sub-problem identified,
not yet resolved within the FST framework. The Yang--Mills line
is included to mark the analogue: UCU3 (global) for the YM functional
\emph{is} the mass-gap question itself.}
\label{tab:status}
\end{table}

\paragraph{Reading the table.}
For thermodynamics, chemistry and biology, UCU1+UCU2 are obtained
from convexity of free-energy-like functionals on physically motivated
feasibility sets; UCU3 is local but not global, because non-equilibrium
boundary conditions can produce competing basins. For K41 turbulence,
the local UCU3 is achieved at the Kolmogorov scaling exponent
$g^*=\tfrac12$, but a logarithmic Sobolev inequality (LSI) is required
to upgrade to a global statement. For dark energy (CRM-VI), UCU1--UCU3
hold for the de~Sitter functional $\Phi(\rho)$ on $\mathbb{R}_+$,
but lifting the result through the L4 RG-matching to a full effective
field theory is not yet completed.

\subsection{Status of FST Building Blocks (mathematics and physics)}
\label{subsec:building-blocks}

Beyond the UCU-application table above, the FST programme includes a
set of published \emph{building blocks} in the mathematics and physics
lines. Table~\ref{tab:building-blocks} summarises their honest proof
status (drawn from each project's individual proof note, see the
companion repository at \texttt{STATUS\_UEBERSICHT.md}).

\begin{table}[h]
\centering
\renewcommand{\arraystretch}{1.18}
\small
\begin{tabular}{p{2.8cm}p{2.0cm}p{2.5cm}p{4.3cm}}
\toprule
\textbf{Building block} & \textbf{Line} & \textbf{Status} &
\textbf{Open core gap} \\
\midrule
RH v2.1 (Even Dominance) & mathematics & unconditional & analytic
$c_{\text{gap}}$ derivation removes last numerical input \\
Zeta Zoo / NE-B          & mathematics & conditional   & Prime-Hub /
OP5 obstructions \\
Hodge                    & mathematics & framework     & Hard direction
(AP $\Rightarrow$ algebraic) \\
BSD                      & mathematics & framework     & Higher
Gross--Zagier for rank $\ge 2$ \\
P vs NP (Entropy)        & mathematics & reformulation & Uniformity
bridge (NOT a proof) \\
\midrule
Yang--Mills mass gap     & physics     & strong-coupling proven &
RG/continuum transport conditional \\
Navier--Stokes (Attractor) & physics   & conditional framework &
Attractor regularity package, Assumption G/G' \\
NS-LDI                   & physics     & framework + proof of life &
TLL+LDI for 3D NS attractors \\
K41 Turbulence (Paper A) & physics     & unconditional & --- (Paper B
remains conditional on $\mathrm{DFC1}^{\vee}$) \\
\midrule
FST-DE (Dark Energy)     & cosmology   & framework + no-go &
UV renormalisation, screening lift \\
\bottomrule
\end{tabular}
\caption{Honest proof status of FST building blocks in the
mathematics, physics and cosmology lines, as of 2026-05. Every entry
references an individual project whose primary proof note (BEWEISNOTIZ.md)
documents the precise lemma chain and the residual open gaps. The
present overview neither proves nor enlarges any of these results;
it locates each within the same UCU/RG/MEPP architecture. Where the
status reads ``reformulation'' (P vs NP) or ``framework'' (Hodge,
BSD, NS, FST-DE), the open core gap is explicitly the original
Millennium-class difficulty -- consistent with the Pattern-A lesson
(Section~\ref{sec:pattern-a}).}
\label{tab:building-blocks}
\end{table}

% ===================================================================
\section{The Pattern-A Lesson}
\label{sec:pattern-a}
% ===================================================================

Across both the mathematics line (Hodge, Beal/abc, BSD, P~vs~NP) and
the cosmology line (CRM-VI, dark energy), one methodological pattern
has recurred so often that we treat it as a programme-level
\emph{warning}.

\begin{quote}
\textbf{Pattern A (negative form).}
A reformulation of an open conjecture $V$ into a ``positivity
normal form'' $P$ is, in general, only a translation of language:
typically $P \Leftrightarrow V$. Without an independent strength
comparison, $P$ is not a proof of $V$.
\end{quote}

\noindent The cosmological variant is operationally identical:
mapping a \emph{physical} question $V_{\text{phys}}$ onto a model
$M$ that contains an open mathematical sub-question $S$ does not
solve $V_{\text{phys}}$; it merely reduces it conditionally to $S$.
The cosmological CC-problem and the inflationary initial-condition
problem each met this pattern in the CRM-VI line. The lesson is
recorded in \texttt{CORE/PATTERN\_A\_LESSON.md} and is the reason
why every UCU-application paper now contains a section labelled
\emph{Honest Status}.

% ===================================================================
\section{Open Sub-Problems}
\label{sec:open}
% ===================================================================

The programme leaves three sharply identified open problems.

\paragraph{HD3' for Beal.}
The Beal-Verifier-Sieve framework reduces the Beal conjecture to four
sub-statements HD3'-A through HD3'-D. As documented in
\texttt{HD3\_B\_EFFECTIVE\_CHARACTER\_SUMS.md} and
\texttt{HD3\_ACD\_TRIO\_OBSTRUCTIONS.md}, the strong character-sum
bound HD3'-B (cancellation $q^{1+\varepsilon}$ rather than the
standard square-root cancellation $q^{3/2}$) lies beyond standard
Weil--Deligne tools, and HD3'-A (Salt-Closing) and HD3'-C (Counting)
reduce to it. A note documenting the conditional structure of HD3'-A
under HD3'-B is provided in \texttt{HD3\_A\_SALT\_CLOSING.md}.

\paragraph{Global UCU3 for Yang--Mills.}
The mass-gap question for $4d$ pure Yang--Mills is exactly the
question whether the energy functional $\mathcal{E}[A]$ admits a
\emph{global} quadratic lower bound around the vacuum: a global UCU3.
This identification means that the FST framework does not solve
Yang--Mills; it locates the mass-gap question as the canonical
UCU3-global representative. UCU1 and UCU2 in the YM context are
themselves open within the FST language.

\paragraph{L4 RG-matching for CRM-VI.}
The cosmological dark-energy line (CRM-VI) supplies UCU1--UCU3 on a
reduced configuration space (de~Sitter functional in
$\mathbb{R}_+$). Lifting the result through a four-step RG-match
(L1 microscopic action, L2 effective potential, L3 cosmological
background, L4 perturbations + matching to data) leaves L4 open.
This is the cosmological analogue of the Pattern-A lesson: the
mathematical reduction is conditional on a physical RG-match that
has not yet been performed at the level required for journal
acceptance.

% ===================================================================
\section{Honest Status}
\label{sec:honest-status}
% ===================================================================

\begin{description}[leftmargin=2em,style=nextline]
\item[What FST delivers.]
A common architecture (universal game $+$ MEPP axiom $+$ UCU
three-lemma endgame $+$ RG bridge) under which scales as different
as the Standard-Model parameter space and dark-energy late-time
acceleration admit a single conditional uniqueness/stability
statement. A status map (Table~\ref{tab:status}) showing where
the architecture is closed and where it is conditional. A
methodological lesson (Pattern~A) which has already prevented
multiple false-progress claims in the mathematics line.

\item[What FST does \emph{not} deliver.]
A solution of any individual hard open problem (Yang--Mills mass
gap, Beal, BSD). FST fixes the \emph{shape} of the conditional
statement; the application-specific UCU3-global step is, in each
hard case, identified with the original difficulty. The novelty
is therefore architectural and reductive, not constructive.

\item[Falsification criterion.]
The programme is falsified if, for any of the listed scales, UCU1
or UCU2 fails on the natural feasibility set. UCU3-global failure
is not a falsification: it is the locator of the open sub-problem
at that scale.
\end{description}

% ===================================================================
\section*{Acknowledgements}
% ===================================================================
This article is a programme overview. The author thanks collaborators
(human and AI) involved in the FST-I, FST-II, FST-III and CRM lines
of the programme. Computations supporting the cosmological line
(Watkins-style spectrum experiments) were performed on
\texttt{ellmos-services}.

% ===================================================================
\section*{AI Disclosure (Level L3)}
% ===================================================================
This paper was prepared with substantial AI assistance (Anthropic
Claude, OpenAI GPT, Google Gemini variants). AI tools were used for:
cross-paper synthesis between the FST-I, FST-II, FST-III, FST-DE
companions; structural editing and consistency checking of the
UCU/RG architecture; generation and verification of the cross-scale
status maps (Tables~\ref{tab:series-components},
\ref{tab:status} and~\ref{tab:building-blocks}); literature
cross-referencing with Zenodo DOIs; and bilingual EN/DE
synchronisation. All scientific claims, the UCU three-lemma
architecture, the Pattern~A lesson, and the falsification criterion
are the author's own and are subject to standard authorial
responsibility. No section was written autonomously by AI.
The categorisation follows the FST programme's five-level AI
disclosure scheme (L1 = no AI; L5 = AI as principal contributor);
this paper is L3 (AI as substantial assistant, human as principal
author and responsible scientist).

% ===================================================================
% Bibliography (skeleton -- to be completed in v1)
% ===================================================================
\begin{thebibliography}{99}

\bibitem{MartyushevSeleznev2006}
L.M.~Martyushev, V.D.~Seleznev,
\emph{Maximum entropy production principle in physics, chemistry and biology},
Phys.~Rep.~\textbf{426} (2006), 1--45.

\bibitem{Brezis2011}
H.~Brezis,
\emph{Functional Analysis, Sobolev Spaces and Partial Differential Equations},
Springer, 2011.

\bibitem{EkelandTemam1999}
I.~Ekeland, R.~T\'emam,
\emph{Convex Analysis and Variational Problems},
SIAM Classics, 1999.

\bibitem{Evans2010}
L.C.~Evans,
\emph{Partial Differential Equations}, 2nd~ed.,
Amer.~Math.~Soc., 2010.

\bibitem{Wilson1974}
K.G.~Wilson, J.~Kogut,
\emph{The renormalisation group and the $\varepsilon$-expansion},
Phys.~Rep.~\textbf{12} (1974), 75--199.

\bibitem{Jacobson1995}
T.~Jacobson,
\emph{Thermodynamics of spacetime: the Einstein equation of state},
Phys.~Rev.~Lett.~\textbf{75} (1995), 1260--1263.

\bibitem{MaynardSmith1982}
J.~Maynard~Smith,
\emph{Evolution and the Theory of Games},
Cambridge Univ.~Press, 1982.

% --- FST programme companion papers (hasPart) ----------------------
\bibitem{GeigerFST-I-2026}
L.~Geiger,
\emph{Functional Stability Theory I: A Game-Theoretic Framework
for the Thermodynamic Stability of Fundamental Parameters},
Zenodo: \texttt{10.5281/zenodo.20130544}, 2026.

\bibitem{GeigerFST-II-2026}
L.~Geiger,
\emph{Functional Stability Theory II: Chemical Stability and
Autocatalytic Selection},
Zenodo: \texttt{10.5281/zenodo.20130563}, 2026.

\bibitem{GeigerFST-III-2026}
L.~Geiger,
\emph{Functional Stability Theory III: Biological Stability and
Nash Frustration},
Zenodo: \texttt{10.5281/zenodo.20130573}, 2026.

\bibitem{GeigerFST-DE-2026}
L.~Geiger,
\emph{Dark Energy as Residual Vacuum Free Energy: a Convexity-Screening
No-Go Theorem within the FST Programme},
Zenodo: \texttt{10.5281/zenodo.19036235}, 2026.

% --- FST mathematical building blocks (references) -----------------
\bibitem{GeigerRH-EvenDom-2026}
L.~Geiger,
\emph{The Riemann Hypothesis via Even Dominance (Trilogy, v2.1)},
Zenodo: \texttt{10.5281/zenodo.19546593}, 2026.

\bibitem{GeigerZetaZoo-2026}
L.~Geiger,
\emph{Zeta Zoo: An NE-B Boundary Classification},
Zenodo: \texttt{10.5281/zenodo.19673227}, 2026.

\bibitem{GeigerZookeeper-2026}
L.~Geiger,
\emph{The Resolvent-Stability Reformulation of the Riemann Hypothesis
(CCM/Zookeeper Route)},
Zenodo: \texttt{10.5281/zenodo.19673127}, 2026.

\bibitem{GeigerSpectrumDuality-2026}
L.~Geiger,
\emph{FST Spectrum Duality (v0.8)},
Zenodo: \texttt{10.5281/zenodo.19672198}, 2026.

\bibitem{GeigerHodge-2026}
L.~Geiger,
\emph{Hodge Conjecture as an Arithmetic Positivity Statement},
Zenodo: \texttt{10.5281/zenodo.20109058}, 2026.

\bibitem{GeigerBSD-2026}
L.~Geiger,
\emph{BSD Conjecture and Néron--Tate Height Positivity},
Zenodo: \texttt{10.5281/zenodo.19087443}, 2026.

\bibitem{GeigerPvsNP-2026}
L.~Geiger,
\emph{P vs NP via an Entropic No-Go Theorem},
Zenodo: \texttt{10.5281/zenodo.19056809}, 2026.

% --- FST physical building blocks (references) ---------------------
\bibitem{GeigerYM-2026}
L.~Geiger,
\emph{Functional Stability Theory: Volume-Independent Spectral Gap
for Lattice Yang--Mills},
Zenodo: \texttt{10.5281/zenodo.19087433}, 2026.

\bibitem{GeigerNS-2026}
L.~Geiger,
\emph{Navier--Stokes Regularity via an Attractor-Distance
Stability Argument},
Zenodo: \texttt{10.5281/zenodo.19087449} (Concept), 2026.

\bibitem{GeigerNSLDI-2026}
L.~Geiger,
\emph{Log-Distance Integrability for Navier--Stokes: a TLL/LDI
Bridge Framework},
Zenodo: \texttt{10.5281/zenodo.19056807}, 2026.

\bibitem{GeigerK41-2026}
L.~Geiger,
\emph{K41 as a Joint Minimiser of a Scale-Resolved Free-Energy
Functional},
Zenodo: \texttt{10.5281/zenodo.19056813}, 2026.

% --- Older series and internal notes -------------------------------
\bibitem{GeigerCRM2026}
L.~Geiger,
\emph{Cosmological Rotation Model -- Papers I--IV},
Zenodo: \texttt{10.5281/zenodo.18728936}, 2026.

\bibitem{FST-Unified2026}
L.~Geiger,
\emph{Pattern A as the Recurring Stability Logic across the FST
Research Ecosystem}, internal note (CORE/PATTERN\_A\_LESSON.md),
2026.

\end{thebibliography}

\end{document}
